\chapter{The Full Multiscale Graph VFE}
\label{ch:full-graph-meta-agent-vfe}

The exact variational free energy was already defined in \Cref{ch:elbo},
extended to one normalized interacting population in
\Cref{ch:local-collective-elbo}, and transported through one declared coarse
channel in \Cref{ch:agent-network-rg}. This chapter does not replace those
results. It asks what additional probabilistic structure is required to put
an adaptive two-channel graph, emergent partitions, and parent-to-child
influence inside one multiscale VFE.

The answer has two levels. The exact answer is always the VFE of one
normalized generative law and one normalized recognition law. A sum of self,
edge, partition, and cross-scale terms is exact only when it is obtained by
factorizing those two laws. Without such a factorization the same sum can be a
useful composite descent potential, but it is not automatically an ELBO and
does not inherit an evidence interpretation. \status{ESTABLISHED}

\section{The full answer before decomposition}
\label{sec:full-graph-vfe-global}

Fix once and for all one base point $c_*\in\mathcal C$ and write
$X_*:=X(c_*)$ for its structural data. Every graph variable, belief or model
presentation, transport, membership kernel, and retained mark in this chapter
is evaluated in the fiber over $c_*$. We suppress that fixed argument. The
symbol $\ell$ denotes RG depth, while $f_\ell$ is reserved for an actual base
coarse map; no such map is present in this fixed-base construction. Fix an
admitted observation $o\in\mathsf O_{\theta,X_*}^{\mathrm{reg}}$ and a finite
depth $L$. Let $W$ contain every random object retained by the multiscale
model: all belief and model coordinates at depths $0,\ldots,L$, edge and
connection variables, partition or membership variables, retained holonomy
marks, and auxiliary latents. Declare one normalized generative kernel
$\mathbb P_\theta(do,dW\given X_*)$ before recognition and one normalized
recognition kernel $\mathbb Q_\phi(dW\given o,X_*)$. With selected posterior
$\boldsymbol\Pi_{\theta,o,X_*}$ and evidence density
$p_\theta(o\given X_*)\in(0,\infty)$, the full multiscale VFE is
\begin{equation}
\boxed{
\mathcal F_{\mathrm{tower}}^{\mathrm{ext}}
[\mathbb Q_\phi;X_*,o]
=-\log p_\theta(o\given X_*)
+\KL\left(
\mathbb Q_\phi(\cdot\given o,X_*)
\middle\Vert
\boldsymbol\Pi_{\theta,o,X_*}
\right).}
\label{eq:full-graph-vfe-global}
\end{equation}
Equivalently, on the finite-density domain,
\begin{equation}
\mathcal F_{\mathrm{tower}}
=\E_{\mathbb Q_\phi}
\left[
\log\frac{q_\phi(W\given o,X_*)}{p_\theta(o,W\given X_*)}
\right].
\label{eq:full-graph-vfe-density}
\end{equation}
Equations~\eqref{eq:full-graph-vfe-global} and
\eqref{eq:full-graph-vfe-density} are the complete VFE. Every later term is an
accounting decomposition of this object or is labeled as a proposed composite
potential. \status{ESTABLISHED}

\section{A normalized hierarchical graph model}
\label{sec:full-graph-vfe-joint}

At depth $\ell$, let $V_\ell$ be the vertex set and let
$Z_\ell=(Z_{i,\ell})_{i\in V_\ell}$ contain the full probabilistic coordinates
of that depth, including belief and model presentations rather than only their
marginals. Put
\begin{equation}
\eta_\ell=(\eta^{b,\ell},\eta^{m,\ell}),
\qquad
G_\ell=(\eta_\ell,\Omega^{b,\ell},\Omega^{m,\ell},\mathsf h_\ell),
\label{eq:full-graph-closed-scale-datum}
\end{equation}
where $\mathsf h_\ell$ contains the retained root, based-holonomy, and dressed
boundary marks defined below. The occupancies and normalized attention rows
are obtained from $\eta_\ell$ by marginalization and disintegration, so they
are not substituted for the edge-event laws in the recursive state. Let
$R_\ell$ be a possibly stochastic membership assignment from $V_\ell$ to
$V_{\ell+1}$. The fixed structural configuration $X_*$ and observation $O$
stay outside the coarse channels.

For adaptive structural coarsening, set $\mathsf Y_\ell=(Z_\ell,G_\ell)$ and
declare one normalized recognition-independent Markov channel
\begin{equation}
C^{\mathrm{struct},c_*}_{\ell+1\leftarrow\ell}:
\mathsf Y_\ell\rightsquigarrow
\mathsf R_\ell\times\mathsf Y_{\ell+1}.
\label{eq:full-graph-common-structural-channel}
\end{equation}
Its output includes $(R_\ell,\eta_{\ell+1},\mathsf h_{\ell+1})$. This channel
supplies the conditional law of $(R_\ell,\mathsf Y_{\ell+1})$ given
$\mathsf Y_\ell$. Define the augmented transition state
\begin{equation}
\widehat{\mathsf Y}_{\ell+1\leftarrow\ell}
:=(\mathsf Y_\ell,R_\ell,\mathsf Y_{\ell+1}).
\label{eq:full-graph-augmented-transition-state}
\end{equation}
Apply the same channel separately to the fine generative marginal, its
selected posterior, and recognition law to obtain the transition joint laws
\begin{align}
\widehat{\mathbb P}^{\mathrm{struct}}_{\ell+1\leftarrow\ell}
(dy_\ell,dr_\ell,dy_{\ell+1})
&:=\mathbb P_\ell(dy_\ell)
C^{\mathrm{struct},c_*}_{\ell+1\leftarrow\ell}
(y_\ell;dr_\ell,dy_{\ell+1}),
\nonumber\\
\widehat{\boldsymbol\Pi}^{\mathrm{struct}}_{\ell+1\leftarrow\ell}
(dy_\ell,dr_\ell,dy_{\ell+1})
&:=\boldsymbol\Pi_\ell(dy_\ell)
C^{\mathrm{struct},c_*}_{\ell+1\leftarrow\ell}
(y_\ell;dr_\ell,dy_{\ell+1}),
\nonumber\\
\widehat{\mathbb Q}^{\mathrm{struct}}_{\ell+1\leftarrow\ell}
(dy_\ell,dr_\ell,dy_{\ell+1})
&:=\mathbb Q_\ell(dy_\ell)
C^{\mathrm{struct},c_*}_{\ell+1\leftarrow\ell}
(y_\ell;dr_\ell,dy_{\ell+1}).
\label{eq:full-graph-common-structural-pushforwards}
\end{align}
These joint laws genuinely retain $(R_\ell,\eta_\ell,\mathsf h_\ell)$; their
next-depth marginals additionally contain
$(\eta_{\ell+1},\mathsf h_{\ell+1})$.
The induced marginals of $R_\ell$ may differ because the three input laws
differ. The structural channel does not. If the channel is represented by an
auxiliary seed, the same measurable map and seed law define all three
pushforwards, and each realized seed gives the same structural map in all three
applications; a law-dependent realized blocking rule is not admitted. The
observation coordinate remains unchanged. \status{HYPOTHESIS}

A normalized top-down generative factorization is
\begin{align}
\mathbb P_\theta(do,dW\given X_*)
={}&L_\theta(do\given Z_0,X_*)P_L(dZ_L,dG_L\given X_*)
\nonumber\\
&\times\prod_{\ell=0}^{L-1}
P_R^\ell(dR_\ell\given Z_{\ell+1},G_{\ell+1},X_*)
P_G^\ell(dG_\ell\given Z_{\ell+1},R_\ell,X_*)
\nonumber\\
&\times\prod_{\ell=0}^{L-1}
K_\downarrow^\ell(dZ_\ell\given Z_{\ell+1},R_\ell,G_\ell,X_*).
\label{eq:full-graph-generative-factorization}
\end{align}
Every factor is a normalized measurable kernel. In particular,
$K_\downarrow^\ell$ is the mathematical location of downward influence: a parent
belief or model changes the conditional prior of its children. The bottom-up
coarse map belongs to recognition or to an induced pushforward; it is not
inserted as a second incompatible generative arrow at the same time slice.
\status{HYPOTHESIS}

For an exact chain-rule decomposition, disintegrate recognition along the same
ordering:
\begin{align}
\mathbb Q_\phi(dW\given o,X_*)
={}&Q_L(dZ_L,dG_L\given o,X_*)
\nonumber\\
&\times\prod_{\ell=0}^{L-1}
Q_R^\ell(dR_\ell\given Z_{\ell+1},G_{\ell+1},o,X_*)
Q_G^\ell(dG_\ell\given Z_{\ell+1},R_\ell,o,X_*)
\nonumber\\
&\times\prod_{\ell=0}^{L-1}
Q_\ell(dZ_\ell\given Z_{\ell+1},R_\ell,G_\ell,o,X_*).
\label{eq:full-graph-recognition-factorization}
\end{align}
This is a disintegration of a correlated law, not a mean-field assumption. A
bottom-up recognition network may parameterize the same joint; the written
ordering is only an accounting choice. \status{HYPOTHESIS}

When the hierarchy is adaptive, the displayed $P_R^\ell,P_G^\ell$ and
$Q_R^\ell,Q_G^\ell$ are selected disintegrations of the respective joint laws
induced by \eqref{eq:full-graph-common-structural-pushforwards}. They are not
independently chosen structural channels.

\theoremheading{Exact hierarchical decomposition}{thm:full-graph-vfe-decomposition}
Assume the normalized factorizations above, common domination, and the
integrability needed for the finite log-density form. Then
\begin{align}
\mathcal F_{\mathrm{tower}}
={}&-\E_{\mathbb Q_\phi}\log L_\theta(o\given Z_0,X_*)
+\KL(Q_L\Vert P_L)
\nonumber\\
&+\sum_{\ell=0}^{L-1}\E_{\mathbb Q_\phi}
\KL\left(Q_R^\ell(\cdot\given Z_{\ell+1},G_{\ell+1},o,X_*)
\middle\Vert P_R^\ell(\cdot\given Z_{\ell+1},G_{\ell+1},X_*)\right)
\nonumber\\
&+\sum_{\ell=0}^{L-1}\E_{\mathbb Q_\phi}
\KL\left(Q_G^\ell(\cdot\given Z_{\ell+1},R_\ell,o,X_*)
\middle\Vert P_G^\ell(\cdot\given Z_{\ell+1},R_\ell,X_*)\right)
\nonumber\\
&+\sum_{\ell=0}^{L-1}\E_{\mathbb Q_\phi}
\KL\left(Q_\ell(\cdot\given Z_{\ell+1},R_\ell,G_\ell,o,X_*)
\middle\Vert K_\downarrow^\ell(\cdot\given Z_{\ell+1},R_\ell,G_\ell,X_*)\right).
\label{eq:full-graph-vfe-factorized}
\end{align}
Each conditional divergence is averaged over its displayed parent variables
under $\mathbb Q_\phi$. These are the observation, top-prior, partition,
graph, and parent--child contributions. \status{ESTABLISHED}

\noindent\emph{Proof.}
Substitute the two factorizations into the Radon--Nikodym ratio in
\eqref{eq:full-graph-vfe-density} and apply the conditional relative-entropy
chain rule. The extended-real version must use the measure-level chain rule
without subtracting infinities. \hfill$\square$

A deterministic same-time parent $Z_{\ell+1}=C_\ell(Z_\ell)$ and an
independently sampled parent that generates $Z_\ell$ are not two free
descriptions of one joint. One may use the hierarchy above, a deterministic
pushforward, or a normalized undirected compatibility factor. Reciprocal
same-time influence
requires the last option and its partition function. Delayed Wheelerian
feedback instead lets $Z_{\ell+1}(t)$ parameterize the child kernel at
$t+\Delta t$; that is a dynamical extension, not a static ELBO identity.
\status{NOT-CLAIMED}

\section{The scale-local two-channel graph energy}
\label{sec:full-graph-vfe-scale-local}

For arbitrary directed, generally non-flat transports, define
\begin{align}
D_{ij}^{b,\ell}
&:=\KL\left(q_{i,\ell}^b\middle\Vert
(\Omega_{ij}^{b,\ell})_\#q_{j,\ell}^b\right),
&
D_{ij}^{m,\ell}
&:=\KL\left(q_{i,\ell}^m\middle\Vert
(\Omega_{ij}^{m,\ell})_\#q_{j,\ell}^m\right).
\label{eq:full-graph-transported-divergences}
\end{align}
Here and below the suppressed base argument is $c_*$. The edge-event laws in
$G_\ell$ determine their receiver occupancies and normalized source rows by
\begin{equation}
\alpha_i^{x,\ell}:=\sum_j\eta_{ij}^{x,\ell},
\qquad
\beta_{ij}^\ell:=\frac{\eta_{ij}^{b,\ell}}{\alpha_i^{b,\ell}},
\qquad
\gamma_{ij}^\ell:=\frac{\eta_{ij}^{m,\ell}}{\alpha_i^{m,\ell}},
\label{eq:full-graph-edge-event-disintegration}
\end{equation}
for $x\in\{b,m\}$ and positive occupancy. On a zero-occupancy row, choose a
declared normalized version; it does not change the joint edge-event law.
The correct energy-unit row free energies are
\begin{align}
\Phi_{i,\ell}^b
&=\sum_j\beta_{ij}^\ell D_{ij}^{b,\ell}
+\tau_{i,\ell}^b\KL(\beta_i^\ell\Vert\pi_i^{b,\ell}),
\nonumber\\
\Phi_{i,\ell}^m
&=\sum_j\gamma_{ij}^\ell D_{ij}^{m,\ell}
+\tau_{i,\ell}^m\KL(\gamma_i^\ell\Vert\pi_i^{m,\ell}).
\label{eq:full-graph-row-free-energies}
\end{align}
Their conditional minimizers are
\begin{equation}
\beta_{ij}^{\ell,*}
=\frac{\pi_{ij}^{b,\ell}e^{-D_{ij}^{b,\ell}/\tau_{i,\ell}^b}}
{\sum_k\pi_{ik}^{b,\ell}e^{-D_{ik}^{b,\ell}/\tau_{i,\ell}^b}},
\qquad
\gamma_{ij}^{\ell,*}
=\frac{\pi_{ij}^{m,\ell}e^{-D_{ij}^{m,\ell}/\tau_{i,\ell}^m}}
{\sum_k\pi_{ik}^{m,\ell}e^{-D_{ik}^{m,\ell}/\tau_{i,\ell}^m}}.
\label{eq:full-graph-gibbs-rows}
\end{equation}
Thus $\beta_{ij}$ and $\gamma_{ij}$ coevolve with the transported divergences.
Omitting the row relative entropy changes the objective. \status{ESTABLISHED}

The quantity $1/\beta_{ij}$ is a useful monotone interaction length only in a
fixed row. It is directed, depends on all competing sources through row
normalization, diverges at zero, and need not obey a triangle inequality. At
the Gibbs optimum the source-relative surprisal
\begin{equation}
\lambda_{ij}^{b,\ell}:=-\tau_{i,\ell}^b
\log\frac{\beta_{ij}^{\ell,*}}{\pi_{ij}^{b,\ell}}
=D_{ij}^{b,\ell}+\tau_{i,\ell}^b\log Z_i^{b,\ell}
\label{eq:full-graph-attention-length}
\end{equation}
is the transported mismatch plus a row-dependent constant. Raw conductance,
a symmetrized resistance, or diffusion time is required for a genuine graph
metric. \status{ESTABLISHED}

If a scale-local prior is a normalized Gibbs law containing the self terms and
$\sum_{ij}\eta_{ij}^{b,\ell}D_{ij}^{b,\ell}
+\sum_{ij}\eta_{ij}^{m,\ell}D_{ij}^{m,\ell}$, those terms and the global normalizer
belong to the exact VFE. If they are merely appended to an already fixed
state-level VFE, define instead
\begin{align}
\mathcal G_{\mathrm{composite}}
={}&\mathcal F_{\mathrm{population}}
+\sum_\ell\sum_i\left(
\alpha_i^{b,\ell}\Phi_{i,\ell}^b
+\alpha_i^{m,\ell}\Phi_{i,\ell}^m\right)
\nonumber\\
&+\sum_\ell\mathcal F_{\mathrm{partition}}^\ell
+\sum_\ell\mathcal F_{\mathrm{cross}}^\ell
+\sum_\ell\mathcal R_{\mathrm{connection}}^\ell.
\label{eq:full-graph-composite-potential}
\end{align}
This is a composite multiscale free-energy potential, not automatically an
evidence bound. \status{DEFINITION}

\section{Non-flat edges and what the parent retains}
\label{sec:full-graph-vfe-nonflat}

The edge transports are not assumed flat. Use the rooted construction of
\Cref{sec:rg-gauge-cross-scale}: in every connected block $I$, choose a root
$r_I$ and a rooted spanning tree, with nested choices across admitted depths.
For $x\in\{b,m\}$, let $\tau_{I\leftarrow i}^{x,\ell}$ be the ordered tree
transport from $i$ to $r_I$. Under endpoint frame changes $a_i^b=a_i$ and
$a_i^m=b_i$,
\begin{equation}
\tau_{I\leftarrow i}^{x,\ell\prime}
=(a_{r_I}^x)^{-1}\tau_{I\leftarrow i}^{x,\ell}a_i^x.
\label{eq:full-graph-tree-transport-law}
\end{equation}
The non-tree edges define the based holonomy representation
\begin{equation}
H_I^{x,\ell}:\pi_1(\Gamma_I,r_I)\longrightarrow G,
\qquad
H_I^{x,\ell\prime}(\gamma)
=(a_{r_I}^x)^{-1}H_I^{x,\ell}(\gamma)a_{r_I}^x.
\label{eq:full-graph-based-holonomy}
\end{equation}
Nonidentity based holonomy is allowed. A holonomy-blind parent law
$Q_I^{x,\ell}$ is path independent only on the stabilized sector
\begin{equation}
\bigl(H_I^{x,\ell}(\gamma)\bigr)_\#Q_I^{x,\ell}=Q_I^{x,\ell}
\quad\text{for every retained based loop }\gamma.
\label{eq:full-graph-holonomy-stabilization}
\end{equation}
This is not a flatness condition and does not select $I$.

For a marked boundary edge $e:j\to i$ from block $J$ to block $I$, write
$\Theta_e^{x,\ell}$ for its microscopic channel-$x$ transport. Its endpoint
gauge law and its root-framed dressed boundary mark are
\begin{align}
\Theta_e^{x,\ell\prime}
&=(a_i^x)^{-1}\Theta_e^{x,\ell}a_j^x,
\nonumber\\
V_e^{x,\ell}
&=\tau_{I\leftarrow i}^{x,\ell}\Theta_e^{x,\ell}
\bigl(\tau_{J\leftarrow j}^{x,\ell}\bigr)^{-1},
\qquad
V_e^{x,\ell\prime}
=(a_{r_I}^x)^{-1}V_e^{x,\ell}a_{r_J}^x.
\label{eq:full-graph-dressed-edge}
\end{align}
Thus $\mathsf h_\ell$ retains the simultaneous root-framed datum
$(r_I,H_I^{x,\ell},\{V_e^{x,\ell}\})$ and its joint law. On a non-flat graph,
the family of dressed boundary marks need not collapse to one group element.
An exact coarse datum retains their conditional distribution; a reduced mean
must carry a declared residual. Penalizing that residual to zero is an extra
modeling preference, not a consequence of gauge covariance. \status{ESTABLISHED}

\section{Partition variables and graph renormalization}
\label{sec:full-graph-vfe-partitions}

The weights $\beta$ and $\gamma$ do not themselves create a parent random
variable. A hierarchy intended to emerge under the same variational principle
must put $R_\ell$ inside $\mathbb P$ and $\mathbb Q$. A common finite-temperature
form is
\begin{equation}
\mathcal F_{\mathrm{partition}}^\ell
=\E_{Q_R^\ell}[E_{\mathrm{block}}^\ell(R_\ell;Z_\ell,G_\ell)]
+\tau_R^\ell\KL(Q_R^\ell\Vert\Pi_R^\ell),
\label{eq:full-graph-partition-free-energy}
\end{equation}
provided the normalized Gibbs factor represented by this expression is part
of the generative law. The block energy may trade internal transported
mismatch against boundary retention, evaluator closure, and residual cost.
No canonical selector follows from the global VFE alone. \status{HYPOTHESIS}

Coarse attention is obtained by pushing the joint edge-event law, not by
averaging normalized rows. Genuine overlapping assignment requires a
normalized membership Markov kernel
\begin{equation}
m_\ell(I\mid i,R_\ell)\geq0,
\qquad
\sum_{I\in V_{\ell+1}}m_\ell(I\mid i,R_\ell)=1.
\label{eq:full-graph-membership-kernel}
\end{equation}
When the two endpoint assignments are conditionally independent given
$R_\ell$, the coarse directed event law is
\begin{equation}
\eta_{IJ}^{x,\ell+1}
=\sum_{i,j\in V_\ell}\eta_{ij}^{x,\ell}
\E_\mu\left[
m_\ell(I\mid i,R_\ell)m_\ell(J\mid j,R_\ell)
\right],
\qquad x\in\{b,m\},
\label{eq:full-graph-edge-pushforward}
\end{equation}
where $\mu$ is the generative, selected-posterior, or recognition input law
being pushed through the same structural channel. When endpoint assignments
are correlated, replace the product by a normalized joint endpoint-assignment
kernel
\begin{equation}
M_\ell(I,J\mid i,j,R_\ell)\geq0,
\qquad
\sum_{I,J}M_\ell(I,J\mid i,j,R_\ell)=1,
\label{eq:full-graph-joint-endpoint-kernel}
\end{equation}
and use
\begin{equation}
\eta_{IJ}^{x,\ell+1}
=\sum_{i,j\in V_\ell}\eta_{ij}^{x,\ell}
\E_\mu\left[M_\ell(I,J\mid i,j,R_\ell)\right].
\label{eq:full-graph-correlated-edge-pushforward}
\end{equation}
The resulting $\eta_{\ell+1}$ is part of $G_{\ell+1}$, and its occupancies and
rows are recovered by \eqref{eq:full-graph-edge-event-disintegration}. A hard
random partition is the deterministic conditional special case
$m_\ell(I\mid i,R_\ell)=\mathbf 1\{R_\ell(i)=I\}$. Literal unit-weight
membership in several parents has total mass greater than one; it is
finite-measure replication, not a probability pushforward. Arbitrary
overlapping indicator sums therefore have no established probability status.

Nested assignments must satisfy Chapman--Kolmogorov composition. In the
single-endpoint case this means
\begin{equation}
m_{\ell+2\leftarrow\ell}(A\mid i,R_\ell,R_{\ell+1})
=\sum_I m_{\ell+1}(A\mid I,R_{\ell+1})
m_\ell(I\mid i,R_\ell),
\label{eq:full-graph-membership-composition}
\end{equation}
with the analogous normalized composition
\begin{equation}
M_{\ell+2\leftarrow\ell}(A,B\mid i,j,R_\ell,R_{\ell+1})
=\sum_{I,J}M_{\ell+1}(A,B\mid I,J,R_{\ell+1})
M_\ell(I,J\mid i,j,R_\ell).
\label{eq:full-graph-joint-membership-composition}
\end{equation}
when correlated endpoint kernels are used. The complete augmented structural
channels, including retained marks, must likewise obey
\begin{equation}
C^{\mathrm{struct},c_*}_{\ell+2\leftarrow\ell}
=C^{\mathrm{struct},c_*}_{\ell+2\leftarrow\ell+1}
\circ C^{\mathrm{struct},c_*}_{\ell+1\leftarrow\ell},
\label{eq:full-graph-structural-channel-composition}
\end{equation}
where the composition integrates all intermediate assignments, edge-event
laws, and marks. Retained root, holonomy, and boundary marks must be included
in the target of each composed channel rather than silently discarded.
\status{ESTABLISHED}

\section{How a parent influences its children}
\label{sec:full-graph-vfe-downward}

The normalized kernel $K_\downarrow^\ell$ supplies the cleanest static answer.
For a block $I$ with children $i\in I$, a factorization such as
\begin{equation}
K_\downarrow^\ell(dZ_\ell\given Z_{\ell+1},R_\ell,G_\ell,X_*)
=\prod_{I\in V_{\ell+1}}K_I^\ell
\left(dZ_I^\ell\given Z_{I,\ell+1},R_\ell,G_\ell,X_*\right)
\label{eq:full-graph-downward-kernel}
\end{equation}
makes the parent belief and model presentations arguments of the children's
conditional generative prior. The corresponding term is
\begin{equation}
\mathcal F_{\mathrm{cross}}^\ell
=\E_{\mathbb Q_\phi}
\KL\left(
Q_\ell(dZ_\ell\given Z_{\ell+1},R_\ell,G_\ell,o,X_*)
\middle\Vert
K_\downarrow^\ell(dZ_\ell\given Z_{\ell+1},R_\ell,G_\ell,X_*)
\right).
\label{eq:full-graph-cross-scale-kl}
\end{equation}
If $K_\downarrow^\ell$ depends nontrivially on the parent coordinates and the
parameterization transmits that dependence to the objective, variation with
respect to a parent coordinate changes the parent update while variation with
respect to a child coordinate changes the child update. Under those
hypotheses, optimization is bidirectionally coupled even though the normalized
generative arrow is top-down. A constant parent-independent downward kernel
decouples this cross-scale term. This is the precise static sense in which a
meta-agent can influence its sub-agents without allowing a generative factor
to read the posterior. \status{ESTABLISHED}

Calling the optimization parameter physical time requires more. Delayed
feedback may use the inferred parent at time $t$ to modify
$K_\downarrow^\ell$ at $t+\Delta t$. A same-time reciprocal factor may instead
be placed in a normalized Gibbs joint. Either construction needs work or flux
accounting before it supports Wheelerian nonequilibrium. Static VFE descent
alone proves neither sustained nonequilibrium nor downward causation in
physical time. \status{OPEN}

\section{Bayesian and network renormalization}
\label{sec:full-graph-vfe-literature}

Bayesian renormalization uses Fisher distinguishability on model-parameter
space to define an information scale and separate stiff from sloppy
directions \citep{BermanKlingerStapleton2023}. It can rank which directions
of a parent model presentation remain empirically distinguishable at a
declared observational precision. It does not by itself select graph blocks,
renormalize non-flat transports, or prove closure of the two-channel
interaction law. Its parameter-space orientation must not be identified
silently with sample-space graph diffusion.

Network renormalization addresses the absence of a lattice, momentum, and
geometric length in a heterogeneous graph \citep{gabrielli2025network}.
Laplacian RG offers a diffusion scale and spectral diagnostics
\citep{villegas2023laplacian}; the multiscale model gives a geometry-free
closure family based on additive hidden variables
\citep{Garuccio2023Multiscale}. For the present theory these are candidate
partition and closure mechanisms. They do not automatically handle directed
row-normalized $\beta,\gamma$, two gauge representations, non-flat holonomy,
or simultaneous renormalization of beliefs, models, topology, and dynamics.
The network-RG review identifies the last problem as open.

The proposed synthesis is hybrid: use edge-event laws or raw conductances to
construct a gauge-invariant multiplex graph operator; use spectral or
Bayesian information diagnostics to propose scales and relevant directions;
infer $R_\ell$ through a normalized partition factor; push the full
probabilistic law and edge-event measures; retain non-flat holonomy or prove
stabilization; and measure failure of the parent family by an explicit closure
residual. \status{HYPOTHESIS}

\section{The next theorem}
\label{sec:full-graph-vfe-target}

Starting from the fixed $c_*\in\mathcal C$ with a finite directed multiplex
datum $G_0=(\eta_0^b,\eta_0^m,\Omega_0^b,\Omega_0^m,\mathsf h_0)$, construct a
normalized partition law and parent--child kernels such that descent of
\eqref{eq:full-graph-vfe-global} produces persistent nested blocks, the
pushforward rules \eqref{eq:full-graph-edge-pushforward}--
\eqref{eq:full-graph-structural-channel-composition} compose across depths,
the parent family has a controlled closure residual, and non-flat holonomy is
retained or stabilized without imposing flatness. Then determine whether the
coupled variational dynamics has a separated slow manifold on which
depth-$\ell+1$ variables feed back consistently into depth-$\ell$ kernels.

No current result proves that arbitrary descent from arbitrary
$\beta,\gamma$ produces such a hierarchy. The exact VFE tells us how a
proposed hierarchy must be normalized and accounted for; it does not supply
the partition selector, persistence theorem, or dynamical cross-scale
closure. Those are the next proof obligations. \status{OPEN}
